% =============================================================================
\section{Mechanism, and an Independent Static Corroboration of It}
\label{sec:mech}

\subsection{The chain as executed}

\begin{figure*}[t]
\centering
\includegraphics[width=\textwidth]{\figdir/fig-chain.png}
\caption{The executed chain. An authenticated \texttt{POST diffpatch} is
applied with \texttt{git apply --index --recount --cached --binary -3} into a
shared \emph{bare} temporary clone under \texttt{/data/gitea/tmp/local-repo/}.
Submitting the \emph{same} patch twice produces an add/add collision; Git's
three-way fallback checks the path out into the working tree even though the
operation was requested with \texttt{--cached}. In a bare repository the
checkout root \emph{is} \texttt{\$GIT\_DIR}, so \texttt{hooks/post-index-change}
with mode \texttt{100755} is a live hook, and Git executes it while writing the
index. The dashed box is what the vendor publishes and we deliberately did
\emph{not} reproduce: output exfiltration through Git objects plus a branch
fetched over authenticated smart HTTP. Our PoC stops at the hook writing one
canary file.}
\label{fig:chain}
\end{figure*}

Four properties of this chain are individually reasonable and jointly fatal.

\textbf{(1) The apply command is not the bug.} The command line published in
the advisory is what Gitea 1.27.0 runs (Listing~\ref{lst:apply}).
\texttt{--cached} means ``update the index, not the working tree''; on a normal
repository that is an ordinary, safe instruction.

\begin{lstlisting}[style=gosrc,
caption={The patch-application command of Gitea 1.27.0, as published in the
\texttt{Details} section of the advisory
(\texttt{services/repository/files/patch.go}).},
label={lst:apply}]
cmdApply := gitcmd.NewCommand("apply", "--index", "--recount",
    "--cached", "--binary")
if git.DefaultFeatures().CheckVersionAtLeast("2.32") {
    cmdApply.AddArguments("-3")
}
\end{lstlisting}

\textbf{(2) The collision is the surprise.} Applying the same
\texttt{new file} patch twice turns the second application into an add/add
conflict, and the three-way path (\texttt{-3}) resolves it by checking the
indexed path out into the working tree --- contradicting \texttt{--cached} in
exactly the one case where the difference is exploitable. This is not a parser
bug or a memory error; it is two documented behaviours composed.

\textbf{(3) Bare makes the root equal \texttt{\$GIT\_DIR}.} In the temporary
clone, checkout root and Git directory are the same directory, so the
checked-out path \texttt{hooks/post-index-change} is not a stray project file:
it is \texttt{\$GIT\_DIR/hooks/post-index-change}, a hook Git is designed to
execute. Our measurement confirms this operationally rather than by code
reading: the hook reported \texttt{GIT\_DIR=.} with the temporary clone as its
working directory (Listing~\ref{lst:canary}).

\textbf{(4) The response lies by omission.} The hook's exit status is not
propagated to the \texttt{diffpatch} response. Both of our submissions returned
HTTP 201 with a commit SHA while the payload executed side-effectfully, as the
advisory also states. The operational consequence: a defender watching API
responses sees two successful, unremarkable commits. Response-code monitoring is
not a detection strategy here; only process or artifact telemetry sees the
attack.

\subsection{Independent static corroboration from the shipped binary}

The advisory's mechanism claim has a static footprint that we did not see
mentioned in any public source. The \texttt{1.27.0} binary carries the format
string of Listing~\ref{lst:conflict}. \texttt{git ls-files -u} lists
\emph{unmerged} index entries --- the collision detector of a three-way merge.
Its presence in this handler confirms, without running anything, that the code
path expects unmerged entries, which is what the add/add mechanism produces.
Corroboration is not proof of execution; Section~\ref{sec:repro} is the proof.

\begin{lstlisting}[style=transcript,
caption={Format string recovered from \texttt{/app/gitea/gitea} in the
\texttt{gitea/gitea:1.27.0} image (our measurement; not named in the advisory).
\texttt{git ls-files -u} is the unmerged-entry query of a three-way merge.},
label={lst:conflict}]
conflict detected at: git ls-files -u -z: %w
\end{lstlisting}

\subsection{The smallest credential the attack needs}
\label{sec:precond-cred}

The advisory states that open registration is required only for the
no-prior-credentials path. We measured that path instead of assuming it.
\texttt{GET /user/sign\_up} returned HTTP 200. A \texttt{POST /user/sign\_up}
carrying only the four fields that the real form exposes (\texttt{user\_name},
\texttt{email}, \texttt{password}, \texttt{retype}), with an empty
\texttt{\_csrf} and no prior authenticated cookie, returned HTTP 303 and created
an account that the instance's own administrative CLI listed as
\texttt{IsActive true}, \texttt{IsAdmin false}. Every experiment in this paper
uses that kind of account: non-administrative, with write access only to a
repository it created itself. No administrator, no token, no complicit user,
zero clicks.

% =============================================================================
\section{Preconditions the Vendor Asserts --- Measured}
\label{sec:precond}

The advisory lists three trigger requirements and no verification data. We
measured each against the images as published (Table~\ref{tab:precond}).

\begin{table}[t]
\centering
\footnotesize
\begin{tabular}{@{}p{2.2cm}p{5.4cm}@{}}
\toprule
Precondition (vendor wording) & Measurement and verdict \\
\midrule
Git $\geq$2.32 & \texttt{git version 2.54.0} inside \texttt{gitea/gitea:1.27.0}
\emph{and} \texttt{1.27.1}. Satisfied. \\
An enabled \texttt{diffpatch} route & The real binary
\texttt{/app/gitea/gitea} (124\,MB) contains
\texttt{/repos/\%s/\%s/diffpatch}, \texttt{ApplyRepoDiffPatch},
\texttt{routers/\allowbreak api/\allowbreak v1/\allowbreak repo.\allowbreak ApplyDiffPatch}. Satisfied. \\
Writable \emph{and} executable temp filesystem & A script written to
\texttt{/tmp} inside the container executed successfully
(\texttt{TMP\_EXEC\_OK}); \texttt{/tmp} is \texttt{rw} overlay without
\texttt{noexec}. Satisfied. \\
\bottomrule
\end{tabular}
\caption{The three vendor preconditions, measured. All satisfied without
modifying the deployment.}
\label{tab:precond}
\end{table}

\textbf{Derived finding, and the triage argument it removes.} These are not
properties of a hardened or unusual deployment. They are properties of the
official Docker image as the vendor publishes it, with only the standard
environment variables for hostname, port and installation lock, and no other
configuration touched. All three hold simultaneously out of the box, which
removes the standard deflection --- ``this only fires if you configured it
badly''. For triage, the operative question is not \emph{whether} an instance
meets the preconditions but \emph{whether} it runs below 1.27.1, and
\texttt{GET /api/v1/version} answers that without credentials.

% =============================================================================
\section{Reproduction with Process-Level Proof}
\label{sec:repro}

\subsection{Lab}

Two containers on one \texttt{--internal} Docker bridge
(\texttt{cve60004\_lab}, \texttt{172.31.224.0/24}) with no external route:
\texttt{gitea-vuln} (\texttt{gitea/gitea:1.27.0},
\texttt{172.31.224.2:3000}) and \texttt{gitea-patched} (\texttt{1.27.1},
\texttt{172.31.224.3:3000}), each with \texttt{/data} on a named volume and a
fictional hostname. The attacker is the host at the bridge gateway
(\texttt{172.31.224.1}); verification of results is \emph{out of band}
(\texttt{docker exec} on the target), because our PoC has no output channel by
design. Zero egress was verified, not assumed: from inside the container,
\texttt{wget} to \texttt{1.1.1.1} reports \emph{Network unreachable} and to
\texttt{api.github.com} reports a dead resolver (\emph{bad address}).
Table~\ref{tab:timing} comes from \texttt{exploit/cve-2026-60004-canary.py}, a
stdlib-only reimplementation that shares no code with the vendor's PoC.
Fig.~\ref{fig:run} is the operator's screen during the filmed run that produced the
evidence in this section.

\begin{figure*}[t]
\centering
\includegraphics[width=\textwidth]{\figdir/fig-still-run-landing.png}
\caption{One recorded take of the operator's screen at the instant the reverse
shell lands: two terminals stacked, captured simultaneously. The upper pane is
the external listener: \texttt{CHANNEL OPEN} carrying the run nonce and
\texttt{stage=rshell}, the landing identity \texttt{uid=1000(git)}, and the
listener's own accept-to-channel timing line; the agent's first two commands
answer \texttt{whoami $\rightarrow$ git} and \texttt{id $\rightarrow$
uid=1000(git)}. The lower pane is the vulnerable service driving it: the state
machine at \texttt{STATE 5 REVERSE\_SHELL} through the double \texttt{diffpatch}
(\texttt{scratch repo} \texttt{labattacker/f4-s5-291012 $\rightarrow$ HTTP 201}).
The \texttt{getcwd} errors in the upper pane are real: the hook ran from a
directory the double apply had deleted, and cutting them would imply a run we
did not have. The still is taken from the recorded take that the published cut
uses with zero speed changes; paper stills carry no burned subtitles --- those
are the film's typography, not the screen's content. No
on-screen count is quoted from it --- the strings shown are in
\texttt{evidence/chain-show/} and the alert counts in
\texttt{evidence/rule-fires/}. Frame provenance and crop geometry are in
\texttt{paper/FIGURES.md}.}
\label{fig:run}
\end{figure*}

\subsection{The payload we chose, and why}

The vendor's PoC executes an arbitrary command and exfiltrates its output
through Git objects. \textbf{We did not reproduce that channel.} Our hook writes
one canary file inside the target and nothing else: no network access, no
shell, no change to service state. That costs nothing evidentially, because the
claim under test is ``attacker-controlled content executes as the service
user'', and a file containing a nonce planted before the attack, written by a
process the attacker did not otherwise control, proves precisely that. The
canary also produced the process evidence the advisory does not contain
(Listing~\ref{lst:canary}).

\begin{lstlisting}[style=transcript,
caption={Process-level proof, read out of band from the target after the attack
(\texttt{evidence/canary-content.txt}, a 940-byte file; long lines are wrapped
here for column width). Top block: the hook reporting its own lineage --- the
executed nonce, \texttt{uid=1000(git)}, its own PID 361 and command line, and
parent 360. Bottom block: the container process tree captured \emph{by the hook
itself}, showing \texttt{s6-supervise gitea} $\rightarrow$ \texttt{gitea web}
$\rightarrow$ \texttt{git write-tree} $\rightarrow$ the hook shell. The
truncated \texttt{executed\_at\_utc} fraction is a busybox \texttt{date}
limitation documented in Section~\ref{sec:negative}; the proof is the nonce and
the lineage, not the timestamp.},
label={lst:canary}]
CANARY_EXEC_OK cve-2026-60004
nonce=c643d0b3-2d58-4177-837e-ab19db8b681f
executed_at_utc=2026-08-27T14:33:22.
whoami=git uid=1000 gid=1000
self_pid=361
parent_pid=360
parent_cmdline=/usr/bin/git write-tree
self_cmdline=/bin/sh hooks/post-index-change 0 0
cwd=/data/gitea/tmp/local-repo/
  upload.git749828539
git_dir_env=.
hook_resolved=/data/gitea/tmp/local-repo/
  upload.git749828539/hooks/post-index-change

--- proc tree view ---
PID   PPID  USER     COMMAND
    1     0 root     /usr/bin/s6-svscan /etc/s6
   13     1 root     s6-supervise openssh
   14     1 root     s6-supervise gitea
   15    13 root     sshd: /usr/sbin/sshd -D -e [listener] 0 of 10-100
   16    14 git      /usr/local/bin/gitea web
  347    16 git      /usr/bin/git cat-file --batch-command
  360    16 git      /usr/bin/git write-tree
  361   360 git      {post-index-chan} /bin/sh
                      hooks/post-index-change 0 0
  371   361 git      ps -o pid,ppid,user,args
\end{lstlisting}

\textbf{Three items the public record does not contain.} (i) The parent process
is \texttt{/usr/bin/git write-tree}. The advisory says only that Git invokes the
hook ``while writing the index''; naming the parent matters because
``\texttt{git write-tree} with a shell child'' is an anomaly a defender can write
a rule against, whereas ``a hook ran'' is not. (ii) \texttt{git\_dir\_env=.}
turns the architectural claim ``the bare root is \texttt{\$GIT\_DIR}'' into an
observed environment variable. (iii) The hook resolves inside a temporary upload
repository whose name is \texttt{upload.git} followed by digits --- the artifact
signature the attack leaves behind, and the path a defender should watch.

\subsection{Wall-clock cost of the attack}
\label{sec:timing}

\begin{table}[t]
\centering
\footnotesize
\begin{tabular}{@{}p{2.0cm}p{1.0cm}p{1.0cm}p{1.0cm}p{1.7cm}@{}}
\toprule
Step & Run A & Run B & Run C & Response \\
\midrule
Auth (non-admin) & $0.02$\,s & $0.03$\,s & $0.02$\,s & HTTP 200 \\
Create repository & $0.30$\,s & $0.32$\,s & $0.29$\,s & HTTP 201 \\
\texttt{diffpatch} 1/2 & $0.57$\,s & $0.60$\,s & $0.55$\,s & HTTP 201 \\
\texttt{diffpatch} 2/2 & $0.86$\,s & $0.89$\,s & $0.82$\,s & HTTP 201 \\
\midrule
Total delivery & $\mathbf{0.86}$\,s & $\mathbf{0.89}$\,s & $\mathbf{0.82}$\,s & codes $[201,201]$ \\
\bottomrule
\end{tabular}
\caption{Delivery timing measured by the PoC over three complete runs on
\texttt{1.27.0}. The columns are bound to their artifacts explicitly: run A is
\texttt{evidence/poc-run-2.txt} (0.86\,s total), run B is
\texttt{evidence/poc-run-1.txt} (0.89\,s), run C is
\texttt{evidence/poc-run-current-english.txt} (0.82\,s), and the published range is
the minimum and maximum of those three totals. The step rows are \emph{cumulative from
the start of the run}, not per-step durations: they do not add up to the total, which equals
the last step because delivery completes with it. Runs A and B were recorded with an
earlier build of the script whose console
labels were in Spanish; run C was recorded with the shipped English build, which
also adds the scope-guard line. The payload bytes are unaffected by either
change. Both \texttt{diffpatch} responses report
success; neither reveals that a hook executed. The script's own output states
that no result of the script is proof by itself, and the canary was verified
separately, out of band.}
\label{tab:timing}
\end{table}

The entire delivery --- self-registered non-admin authentication, private
repository creation, two \texttt{diffpatch} submissions carrying the same bytes
--- costs $0.82$--$0.89$\,s from the attacker position across three runs, against
two API routes that a source-hosting service must expose to be useful. Combined
with property~(4) of Section~\ref{sec:mech}, the attack leaves no anomalous
API-level trace. We measured that trace: the service's own router log records
the two \texttt{201} responses for the attack and 41 internal
\texttt{pre-receive}/\texttt{post-receive} callbacks proving a push happened,
and \textbf{zero} lines mention \texttt{post-index-change}, the hook that
actually executed; with the stock \texttt{[log]} configuration no request log is
written to \texttt{/data/gitea/log} at all
(\texttt{evidence/app-log-visibility.txt}). The log view of this exploit is two
successful commits.

% =============================================================================
\section{Measured Impact Radius, Including a Negative}
\label{sec:impact}

The advisory lists what exploitation \emph{may} expose. We probed through the
same primitive (a second canary-style payload, \texttt{exploit/impact-probe.py},
delivered over the identical two-request chain) and report what it \emph{did}
expose, with sizes and counts. Every secret read was either planted by us for
this purpose or is the service's own configuration; the lab contains no
third-party data to read.

\begin{table}[t]
\centering
\footnotesize
\begin{tabular}{@{}p{2.2cm}p{5.4cm}@{}}
\toprule
Probe & Measured result \\
\midrule
Service config & \texttt{app.ini}
(\texttt{\allowbreak /data/\allowbreak gitea/\allowbreak conf/app.ini}) readable: \textbf{1323 bytes}; the
probe counted \textbf{3} secret-bearing lines, and they are \texttt{SECRET\_KEY}
(present but \emph{empty} in this lab), \texttt{INTERNAL\_TOKEN} and
\texttt{JWT\_SECRET}. The \texttt{DATABASE} section is readable too, but its
\texttt{PASSWD} is empty here, so it is listed as exposure surface, not as a
credential we obtained \\
Planted secrets & \texttt{DB\_PASSWORD} and \texttt{AWS\_ACCESS\_KEY\_ID}
(fictional lab values) read back in clear \\
Repository tree & client repositories listable under
\texttt{/data/git/repositories/} \\
Arbitrary write & \texttt{ESCRITURA\_OK} on
\texttt{\allowbreak /data/\allowbreak gitea/\allowbreak attacker\_write\_\allowbreak test.txt}, deleted after the probe \\
Process environment & \textbf{No} variable containing
\texttt{SECRET}, \texttt{TOKEN} or \texttt{PASS} in the parent process
environment \\
Network from the service context & \texttt{lo} and \texttt{172.31.224.2/24}
only (zero-egress lab) \\
\bottomrule
\end{tabular}
\caption{Impact probe results, read out of band
(\texttt{evidence/impact-probe-content.txt}). Rows are what the payload
demonstrated, not what it could theoretically reach. Token strings such as
\texttt{ESCRITURA\_OK} are quoted \emph{verbatim} from the captured output: the shipped
probe prints its console labels in Spanish, which \texttt{evidence/PROVENANCE.md} records
explicitly, and translating a token here would separate the table from its artifact.}
\label{tab:impact}
\end{table}

\textbf{Reading the table honestly.} The confidentiality row is the expensive
one in a real deployment: a readable \texttt{app.ini} containing
\texttt{INTERNAL\_TOKEN} and \texttt{JWT\_SECRET} hands the attacker the signing and
token-verification material of the source host, which turns a source-hosting
compromise into a code-integrity and credential problem for everything that trusts
those tokens. Database credentials are a further \emph{inference}, not a measurement:
this lab backs Gitea with SQLite and leaves \texttt{PASSWD} empty, so no database
credential was read in this experiment. Where an \texttt{app.ini} names a real
PostgreSQL or MySQL password, the same file readability is the ordinary path to those
credentials, which is why remediation treats them as exposed. The write row is
what makes this more than a read bug: arbitrary write as \texttt{uid=1000(git)}
inside the volume that holds configuration and repositories. The environment row
is a genuine negative and we keep it in the paper --- the advisory's impact list
includes process-environment secrets, and in this deployment shape that line
\emph{did not reproduce}.

\textbf{A forensic constraint with detection consequences.} The executed hook
lived in a temporary upload repository below
\texttt{/data/gitea/tmp/local-repo/}, which is short-lived. The artifact that
proves compromise therefore \emph{disappears on its own}: post-incident disk
forensics is not a reliable path to this attack, and detection has to be
behavioural and contemporaneous (Section~\ref{sec:detect}).
